% Part: conditional-logics
% Chapter: minimal-change-semantics
% Section: sphere-models

\documentclass[../../../include/open-logic-section]{subfiles}

\begin{document}

\olfileid{con}{min}{sph}

\olsection{نماذج الكرات}

إحدى طرق تقديم دلالة صورية للشرطيات المخالفة للواقع هي تحويل تحليل لويس
غير الصوري إلى بنية رياضية. وعندئذ تكون الكرات المحيطة بعالم~$w$
مجموعات من العوالم. ولأن الكرات متداخلة، يجب أن تكون مجموعات العوالم
المحيطة بـ$w$ مرتبة ترتيبًا خطيًا بعلاقة المجموعة الجزئية.

\begin{defn}
  \emph{نموذج الكرات} هو الثلاثي~$\mModel{M} = \tuple{W, O, V}$،
  حيث~$W$ مجموعة غير خالية من العوالم،
  و$V\colon \PVar \to \Pow{W}$ تقييم، و
  $O\colon W \to \Pow{\Pow{W}}$ تسند إلى كل عالم~$w$
  \emph{نظام كرات}~$O_w$. ولكل~$w$، تكون~$O_w$ مجموعة من مجموعات
  العوالم، ويجب أن تحقق:
  \begin{enumerate}
  \item تكون~$O_w$ \emph{متمركزة} عند~$w$: $\{w\} \in O_w$.
  \item تكون~$O_w$ \emph{متداخلة}: كلما كان~$S_1$، $S_2 \in O_w$،
    كان~$S_1 \subseteq S_2$ أو~$S_2 \subseteq S_1$، أي إن~$O_w$
    مرتبة خطيًا بـ$\subseteq$.
  \item تكون~$O_w$ مغلقة تحت الاتحادات غير الخالية.
  \item تكون~$O_w$ مغلقة تحت التقاطعات غير الخالية.
  \end{enumerate}
\end{defn}

الفكرة الحدسية وراء~$O_w$ هي أن العوالم «المحيطة» بـ$w$ مقسّمة إلى
طبقات بحسب بعدها عن~$w$. وأعمق كرة هي~$w$ وحده، أي المجموعة~$\{w\}$:
فـ$w$ أقرب إلى نفسه من العوالم في أي كرة أخرى. وإذا كان
$S \subsetneq S'$، فإن عوالم~$S' \setminus S$ أبعد عن~$w$ من
عوالم~$S$: وتمثل~$S' \setminus S$ «الطبقة» الواقعة بين~$S$ والعوالم
خارج~$S'$. وبوجه خاص، علينا أن نفهم الكرات بوصفها محتوية على جميع
العوالم داخل سطحها الخارجي؛ فهي ليست مجرد الطبقات المفردة.

\begin{figure}
\begin{center}
\begin{tikzpicture}[layerwidth=1.5,scale=.6]\tiny
  \spheresystem{3}
  \propositionintersect{0}{3}{60}{6}
  \path (0,0) node[world] {$w$};
  \spherepos{-30}{2}{node[world] {$w_2$}}
  \spherepos{220}{2}{node[world] {$w_3$}}
  \spherepos{90}{2}{node[world] {$w_1$}}
  \spherepos[shift={(.1,0)}]{10}{3}{node[world] {$w_5$}}
  \spherepos[shift={(.1,0)}]{-10}{3}{node[world] {$w_6$}}
  \spherepos{225}{3}{node[world] {$w_4$}}
  \spherepos{20}{4}{node[world] {$w_7$}}
  \path (5,0) node[left] {$p$};
\end{tikzpicture}
\caption{رسم لنموذج كرات}
\ollabel{fig:sphere-model}
\end{center}
\end{figure}

يقابل الرسم في~\olref{fig:sphere-model} نموذج الكرات الذي فيه
$W = \{w, w_1, \dots, w_7\}$ و$V(p) = \{w_5, w_6, w_7\}$.
أعمق كرة هي~$S_1 = \{w\}$. وأقرب العوالم إلى~$w$ هي~$w_1$ و$w_2$
و$w_3$، ولذلك تكون الكرة الأكبر التالية
$S_2 = \{w, w_1, w_2, w_3\}$. والعوالم الأبعد هي~$w_4$ و$w_5$
و$w_6$، ولذلك تكون الكرة الخارجية
$S_3 = \{w, w_1, \dots, w_6\}$. ونظام الكرات المحيطة بـ$w$ هو
$O_w = \{S_1, S_2, S_3\}$. أما العالم~$w_7$ فلا يقع في أي كرة
محيطة بـ$w$. وأقرب العوالم التي تصدق فيها~$p$ هي~$w_5$ و$w_6$،
ومن ثم تكون أصغر كرة تقبل~$p$ هي~$S_3$.

لتعريف إرضاء صيغة~$!A$ عند عالم~$w$ في نموذج كرات~$\mModel M$،
أي~$\mSat{M}{!A}[w]$، نوسّع تعريف !!{formula} القضايا بإضافة بند
للصيغة~$!B \cif !C$:

\begin{defn}
  $\mSat{M}{!B \cif !C}[w]$ إذا وفقط إذا تحقق أحد الأمرين:
  \begin{enumerate}
  \item\ollabel{sphere-vac} لكل~$u \in \bigcup O_w$،
    $\mSat/{M}{!B}[u]$؛ أو
  \item\ollabel{sphere-nonvac} لبعض~$S \in O_w$،
    \begin{enumerate}
      \item $\mSat{M}{!B}[u]$ لبعض~$u \in S$؛ و
      \item لكل~$v \in S$، إما~$\mSat/{M}{!B}[v]$ أو
        $\mSat{M}{!C}[v]$.
    \end{enumerate}
  \end{enumerate}
\end{defn}

بحسب هذا التعريف، يكون~$\mSat{M}{!B \cif !C}[w]$ إذا وفقط إذا كان
المقدّم~$!B$ كاذبًا في كل موضع من الكرات المحيطة بـ$w$، أو وجدت
كرة~$S$ تصدق فيها~$!B$ ويصدق الشرط المادي~$!B \lif !C$ عند جميع
العوالم في تلك الكرة «التي تقبل~$!B$». لاحظ أننا لم نشترط في التعريف
أن تكون~$S$ \emph{أعمق} كرة تقبل~$!B$، خلافًا لما قد يوحي به الشرح
الحدسي. لكن إذا تحقق الشرط في~\olref{sphere-nonvac} لكرة~$S$، فإنه
يتحقق أيضًا لكل كرة تقبل~$!B$ وتكون محتواة في~$S$؛ ومن ثم، إذا وجدت
أعمق كرة تقبل~$!B$، تحقق الشرط لها على وجه الخصوص.

ولاحظ أيضًا أن تعريف نماذج الكرات لا يشترط أن توجد أعمق كرة تقبل~$!B$:
فقد تكون لدينا متتالية لا نهائية
$S_1 \supsetneq S_2 \supsetneq \dots \supsetneq \{w\}$ من الكرات
التي تقبل~$!B$، ومن ثم لا توجد أعمق كرة تقبل~$!B$. وفي هذه الحالة،
يكون~$\mSat{M}{!B \cif !C}[w]$ إذا وفقط إذا تحقق~$!B \lif !C$
في جميع الكرات~$S_i$، $S_{i+1}$، \dots، لبعض~$i$.

\end{document}
